\subsection{Applications}
\label{sec:applications}

The same momentum and energy calculations are useful in real engineering.
\begin{itemize}
  \item \textit{Attitude control.} The spin from
    Eq.~\eqref{eq:omega-inelastic} is a disturbance that the reaction
    wheels or thrusters have to cancel. The angular momentum from a hit,
    $m v_n b$, tells engineers how much the wheels must be able to absorb.
    For example, Sentinel-1A had a sudden change of orientation after its
    2016 hit \citep{krag2017}.
  \item \textit{Keeping the orbit.} Comparing $|\Delta\vect V|$ with
    $v_{\mathrm{orb}}$ tells us whether a hit needs an orbit correction.
    For debris from millimetres to centimetres in size the change is tiny,
    but a fragment of several kilograms would need a correction burn.
  \item \textit{Shielding and panel design.} Equation~\eqref{eq:ediss}
    shows that in an inelastic hit almost all the energy is released at
    one spot. This is why spacecraft use spaced shields (Whipple shields)
    that break the fragment apart before it reaches the main wall, and why
    solar panels are built so that one crater only knocks out part of the
    power.
  \item \textit{The debris problem.} Every hit can create new fragments.
    Working out the energy and momentum of collisions is the first step in
    models of the debris population and of the chain reaction predicted by
    \citet{kessler1978}.
\end{itemize}
